\source{Лекции, с. 34–35, 38–39; Карпов, §6.1.3, §6.2.1; координатные формулы — пояснение составителя}
\begin{definition}[Сферическая укладка]
Укладка на сфере определяется так же, как на плоскости: вершины — точки
сферы, рёбра — простые дуги без посторонних пересечений.
Все грани на сфере равноправны, выделенной внешней грани нет.
\end{definition}
\begin{lemma}[Стереографическая проекция]
Сфера с удалённой точкой гомеоморфна плоскости.
\end{lemma}
\begin{proof}
Для единичной сферы \(x^2+y^2+z^2=1\) и северного полюса \(N=(0,0,1)\)
проекция на плоскость \(z=0\) задаётся формулой
\[
 (x,y,z)\longmapsto\left(\frac{x}{1-z},\frac{y}{1-z}\right).
\]
Её обратное отображение для \(s=u^2+v^2\):
\[
 (u,v)\longmapsto\left(\frac{2u}{s+1},\frac{2v}{s+1},\frac{s-1}{s+1}\right).
\]
Подстановка проверяет, что образы лежат на сфере и отображения взаимно обратны.
Оба непрерывны, поскольку знаменатели не обращаются в нуль на своих областях.
Для любой другой удаляемой точки предварительно повернём сферу.
\end{proof}
\begin{theorem}[Плоскость и сфера]
Граф допускает укладку на плоскости тогда и только тогда,
когда допускает укладку на сфере.
\end{theorem}
\begin{proof}
Обратная стереографическая проекция переносит плоскую укладку на сферу.
Для обратного перехода выбираем на сфере точку, не лежащую на рисунке,
и проектируем из неё. Такая точка есть: конечное объединение
кусочно-гладких дуг и точек не покрывает сферу.
Гомеоморфизм сохраняет инъективность каждой дуги и отсутствие пересечений.
\end{proof}
\begin{theorem}[Выбор внешней грани]
Любую грань данной сферической укладки можно сделать внешней в некоторой
плоской укладке того же графа.
\end{theorem}
\begin{proof}
Выберем полюс проекции внутри нужной грани. При приближении к полюсу
образ уходит на бесконечность, поэтому образ этой грани неограничен.
Остальные грани ограничены: замыкание каждой не содержит выбранного
полюса и является компактным подмножеством области проекции.
Таким образом, выбранная грань становится внешней.
\end{proof}
\begin{cor}[Вершина и ребро на внешней грани]
Для любой выбранной вершины или любого выбранного ребра планарного
графа существует укладка, в которой этот объект лежит на границе внешней грани.
\end{cor}
\begin{proof}
Выбираем инцидентную объекту грань и делаем её внешней по теореме.
Для вершины такая грань есть в любом свободном секторе её окрестности;
для изолированной вершины — в окружающем её дополнении.
Грань, инцидентная внутренней точке ребра, инцидентна всей его дуге,
поскольку внутри ребра нет вершин. Для моста обе стороны принадлежат
одной грани, что не меняет рассуждения.
\end{proof}
